{\footnotesize
\adjustbox{max width=\linewidth}{\begin{tabular}{p{2.4cm}p{7.2cm}p{4.9cm}}
\toprule
\textbf{Question id} & \textbf{Question (released paraphrase)} & \textbf{Verified answer provision} \\
\midrule
\texttt{ajke\_0144\_02}\newline family\newline colloquial & \bng{৯০ দিন পার হওয়ার পর আমরা যদি আবার একসাথে ঘর-সংসার শুরু করতে চাই, তখন কি হিল্লা বিয়ে করাতেই হবে}? & The Muslim Family Laws Ordinance, 1961, \S 7 (Talaq) [english] \\
\texttt{lawy\_0221\_01}\newline women children\newline colloquial & \bng{নারী ও শিশু নির্যাতন দমন আইনে মামলা হলে জামিন কখন পাওয়া যায়, আর পেতে হলে কী করতে হয়}? & \bng{নারী ও শিশু নির্যাতন দমন আইন, ২০০০}, \S \bng{১৯} [bengali] \\
\texttt{ajke\_0139\_01}\newline consumer\newline colloquial & \bng{আমি ছয় বছর ধরে নিজের নকশায় রেজিনের জিনিস বানিয়ে বিক্রি করছি}। \bng{আমার তৈরি ডিজাইনগুলো কপিরাইটের আওতায় আনতে কী করতে হবে}? & \bng{কপিরাইট আইন, ২০২৩}, \S \bng{৫৩} (\bng{কপিরাইট নিবন্ধন}) [bengali] \\
\texttt{lawy\_0241\_01}\newline land\newline colloquial & \bng{সরকার জমি অধিগ্রহণ করার সময় কোনো বিরোধ বা আপত্তি দেখা দিলে কী পদক্ষেপ নেওয়া উচিত}? & \bng{স্থাবর সম্পত্তি অধিগ্রহণ ও হুকুমদখল আইন, ২০১৭}, \S \bng{৫} (\bng{অধিগ্রহণের বিরুদ্ধে আপত্তি}) [bengali]\newline \bng{স্থাবর সম্পত্তি অধিগ্রহণ ও হুকুমদখল আইন, ২০১৭}, \S \bng{৩০} (\bng{আরবিট্রেটরের নিকট আবেদন}) [bengali] \\
\bottomrule
\end{tabular}}
}
